\section{Closure proofs and exact certificates}\label{cl:appendix}

\subsection{Parameters and coefficient certificates}\label{cl:certificate-tools}

Put
\[
 M_b=M(\sbar),\quad \Xi=F^TM_b,\quad \Sigma=\sym(\Xi),\quad
 N_j=M_b^{-1}M_0(p_j),\quad m=M_b^{-1}e_1,
\]
and $E=e_1e_1^T$, $N_E=M_b^{-1}E$, $Z=\sym(\Xi N_E)$.
Then
\[
 A_F(\sbar+\mu p_j)=\Sigma+\mu\sym(\Xi N_j),\qquad
 v^TZv=(v^T\Xi m)v_1.
\]
A vector $v$ is \emph{kept} if $v^TZv\ge0$. Every $s$ in the upward
completion $B_F$ satisfies $v^TA_F(s)v\ge0$ for every kept $v$:
this follows first for $s=c+te_w$, $c\in C_F$, $t\ge0$, and then by
continuity. All certificate arguments below use this necessary condition.

\begin{lemma}[Transformed point-rule parameters]\label{cl:point-parameters}
The uncompleted point-rule sets under all determinant-preserving
homogeneous automorphisms are exactly those with $\Xi=\Sigma\succ0$
symmetric.
\end{lemma}

\begin{proof}
In coordinates where the chosen point-rule direction is $e_1$, the rule
requires the image matrix to be symmetric with positive trace.
For an automorphism $M\mapsto AMB^T$, $\det A\det B=1$, the preimage of
the positive-semidefinite symmetric-part cone is $C_F$ with
$F^T=B^{-1}A$. At the vertex, the rule gives
$Y=AM_bB^T$ symmetric with positive trace. Since
$\det Y=\det M_b>0$, it is positive definite. Therefore
$F^TM_b=B^{-1}YB^{-T}\succ0$ is symmetric.
Conversely, given $\Sigma\succ0$, set $F^T=\Sigma M_b^{-1}$.
Its determinant is positive. Choose $B=tI$, $A=tF^T$, with
$t^4\det F=1$; then $\det A\det B=1$ and the image of the vertex is
$t^2\Sigma\succ0$. Transposition fixes the base symmetric-part cone
and gives no additional parameter family.
\end{proof}

\begin{lemma}[A parameter half-space for completions]\label{cl:halfspace}
For $\det F>0$ and $\sbar\in\intt B_F$, set $u=(1,-\bar y)^T$.
Then $u^T\Sigma u>0$.
\end{lemma}

\begin{proof}
The affine image of $s\mapsto\sym(F^TM(s))$ has rank three or two.
In the rank-two case a nonzero kernel direction obeys
$M_0(d)=tF^{-T}J$, where $J=\left(\begin{smallmatrix}0&1\\-1&0\end{smallmatrix}\right)$,
so $(F^{-T}J)_{22}=0$. Its image plane cannot contain the apex zero:
$\sym(F^TM(s))=0$ would imply $M(s)=t'F^{-T}J$, contradicting
$M(s)_{22}=1$. A plane meeting the positive-semidefinite cone but missing
its apex cannot support that cone; hence it meets its positive definite
interior. Thus $C_F\ne\varnothing$ implies $\intt C_F\ne\varnothing$.
This also follows from the explicit completion classification in
\cref{fd:completion}.

Interiors of convex sums now give a $\tau>0$ with
$s_\tau=\sbar-\tau e_w\in\intt C_F$.
Let $\Xi_\tau=F^TM(s_\tau)$ and $q_\tau=q(\sbar)-\tau$.
Then $\sym(\Xi_\tau)\succ0$, so $q_\tau>0$, and
\[
 \Xi=\Xi_\tau+\frac{\tau}{q_\tau}\Xi_\tau u e_1^T.
\]
As $e_1^Tu=1$, this gives
$u^T\Sigma u=(1+\tau/q_\tau)u^T\sym(\Xi_\tau)u>0$.
\end{proof}

\begin{lemma}[A parameter-dependent kept vector]\label{cl:kept-boundary}
For a point-rule parameter $\Sigma\succ0$, let $v=J\Sigma m$ and
$L_j(\Sigma)=m^TJN_jJ\Sigma m$. Then
\[
 v^TZv=0,\qquad
 v^T\Sigma v=\det(\Sigma)m^T\Sigma m,\qquad
 -v^T\Sigma N_jv=\det(\Sigma)L_j(\Sigma).
\]
Consequently $a_j(B_F)\ge L_j(\Sigma)/(m^T\Sigma m)$ whenever that ratio
is positive. Each $L_j$ is linear in $\Sigma$.
\end{lemma}

\begin{proof}
The first identity is $(J\Sigma m)^T\Sigma m=0$.
The other two follow from
$J^T\Sigma J=\adj(\Sigma)$,
$\Sigma\adj(\Sigma)\Sigma=\det(\Sigma)\Sigma$, and
$\Sigma J^T\Sigma=-\det(\Sigma)J$.
Applying the kept inequality to $\sbar+\mu p_j\in B_F$ bounds its step
by $m^T\Sigma m/L_j(\Sigma)$.
\end{proof}

\begin{lemma}[Finite box certificate]\label{cl:box-certificate}
Let $\widehat\lambda\ge0$. Suppose a parameter map $\Xi(t)$ is affine
on a box. For distinct selected rays, choose vectors $v_j$ such that at
every box vertex
\[
 v_j^T\Xi N_Ev_j\ge0,\qquad
 n_j=-v_j^T\Xi N_jv_j>0.
\]
Set $\mu_j=\max_{t\text{ vertex}}v_j^T\Xi(t)v_j/n_j(t)$ and assume
$\mu_j>0$. If $\sum_j\widehat\lambda_j/\mu_j\ge1$, then every completion
in the box containing $\sbar$ satisfies
$a^T\widehat\lambda\ge1$.
For point-rule parameters, a positive lower bound from
\cref{cl:kept-boundary} may replace $1/\mu_j$.
Alternatively, a vector kept at all box vertices and satisfying
$v^T\Xi(t)v<0$ at all those vertices excludes every completion in that
box that contains the vertex.
\end{lemma}

\begin{proof}
Keptness and positivity of $n_j$ extend across the box because they
are affine in its parameters. A ratio of affine functions with positive
denominator lies between its values at the vertices: writing a point as
a convex combination weights its vertex ratios by their positive
denominators. Thus for $\mu>\mu_j$,
$v_j^TA_F(\sbar+\mu p_j)v_j<0$ throughout the box.
The kept inequality gives $\alpha_j\le\mu_j$ and then
$a^T\widehat\lambda\ge\sum_j\widehat\lambda_j/\mu_j$.
Unlisted rays contribute a nonnegative amount, and no ray is counted
twice. The exclusion conclusion applies the kept inequality at $\sbar$.
The parameter-dependent ratio is bounded in the same way using its
linear numerator and positive linear denominator.
\end{proof}

\begin{lemma}[A coefficient-simplex certificate]\label{cl:sdp-certificate}
Let $\Pi\subseteq\R_+^N$ be a polytope and let $I_0$ index coordinates
zero on $\Pi$. Suppose symmetric positive-semidefinite matrices $Y_j$
satisfy $Y_j=0$ for $j\in I_0$, and
\[
 R=-\sum_jN_jY_j\quad\text{is symmetric},\qquad
 Q(a)=R-\sum_ja_jY_j\succeq0,\quad\tr Q(a)>0
\]
at every vertex of $\Pi$. No orbit cut vector belongs to
$\Pi+\cone\{e_j:j\in I_0\}$.
\end{lemma}

\begin{proof}
The conditions on $Q$ extend to $\Pi$ by affinity and along its added
coordinates because the corresponding $Y_j$ vanish.
For an orbit set containing the vertex in its interior, $\Sigma\succ0$ and
$a_j\Sigma+\sym(\Xi N_j)\succeq0$; when $a_j=0$ this is the recession
condition. Therefore
\[
 0\le\sum_j\ip{a_j\Sigma+\sym(\Xi N_j)}{Y_j}
 =-\ip{\Sigma}{R-\sum_ja_jY_j}<0,
\]
where the equality uses symmetry of $R$ and the strict inequality uses
$\Sigma\succ0$, $Q\succeq0$, $Q\ne0$.
\end{proof}

To certify a closure point, subdivide
$\{a\ge0:\widehat\lambda^Ta\le1\}$ into simplices and apply
\cref{cl:sdp-certificate} on each. If $\widehat\lambda_j=0$, that
coordinate is an unbounded nonnegative direction and its multiplier must
vanish. A simplex can also be excluded using $x\in X$ with $a^Tx<1$
at each vertex, provided $x_j=0$ on every unbounded coordinate. These
conditions ensure $a^Tx<1$ throughout that simplex and its added cone.

\subsection{Portable finite certificates}\label{cl:finite-certificates}

The electronic supplement contains the exact rational subdivisions and
witnesses in \cref{cl:certificate-list}, together with separate
exact-arithmetic checkers. These finite records establish the stated
membership claims. No numerical search tolerance enters their acceptance
conditions. The theorem statements specify the instance independently of
the records, and the checkers reconstruct $M_b$, $N_j$, $N_E$, and the
parameter charts from the rational instance data.

\begin{table}[htbp]
\centering\small
\caption{Completed exact closure certificates. File names refer to the
electronic supplement.}\label{cl:certificate-list}
\begin{tabular}{llrrr}
\toprule
File & point or conclusion & cut & excluded & skipped\\
\midrule
\texttt{boxcert\_thm14t\_B.json} & $(67/250,41/100,151/500)\in\mathcal C_{\famB}$ &8070&2253&0\\
\texttt{boxcert\_thm14\_B.json} & $(3/11,41/99,10/33)\in\mathcal C_{\famB}$ &2530&725&0\\
\texttt{boxcert\_thm14t\_BP.json} & $(51/350,0,17/150)\in\mathcal C_{\famBP}$ &420&0&6\\
\texttt{boxcert\_thm14\_BP.json} & $(1/5,0,1/6)\in\mathcal C_{\famBP}$ &33&0&5\\
\texttt{boxcert\_thm14w\_B.json} & $(2211,8805,2464)/10000\in\mathcal C_{\famB}$ &5518&1543&0\\
\midrule
\texttt{closure\_cert\_thm14\_A.json} & $(3/11,41/99,10/33)\in\mathcal C_{\famA}$ &25&0&0\\
\texttt{closure\_cert\_prop16\_A.json} & $(24/25,1/100,29/1000)\in\mathcal C_{\famA}$ &92&0&0\\
\texttt{closure\_cert\_wcorner\_A\_A.json} & $(20,20,20,0)\in\mathcal C_{\famA}$ &1&0&0\\
\bottomrule
\end{tabular}
\end{table}

For $\famB$, \cref{cl:halfspace} permits normalization in the seven
facets of $[-1,1]^4$ with first coordinate nonnegative.
More explicitly, $U=[u,e_2]$ and
$U^T\Xi U=\left(\begin{smallmatrix}t_1&t_2\\t_3&t_4\end{smallmatrix}\right)$,
where $t_1>0$. Positive scaling puts one raw entry at $1$ or $-1$,
with the negative first-entry facet omitted. Each chart is therefore an
axis-parallel box of dimension three. The records specify a binary
subdivision along a longest side, choosing the lowest index on ties.
Reconstructing the subdivision tree establishes coverage of each chart;
each terminal box has a cut or exclusion witness from
\cref{cl:box-certificate}. Charts may overlap, which does not affect
coverage or the universal conclusion.

For $\famBP$, normalize the symmetric parameter to trace one:
\[
 \Sigma=\begin{pmatrix}(1+u_1)/2&u_2/2\\u_2/2&(1-u_1)/2\end{pmatrix},
 \qquad u_1^2+u_2^2<1.
\]
The saved certificates use a box containing this disk, with affine
coordinates $u_1=t_1-t_2$, $u_2=t_1+t_2-1$ and
$(t_1,t_2)\in[-1/2,3/2]^2$. A skipped box must have exact minimum
disk polynomial at least one. This minimum is found by evaluating
vertices, edge stationary points, and the interior stationary point.
For the three-ray instance, the parameter-dependent forms are
\[
 m=(2/3,0)^T,\quad m^T\Sigma m=4a/9,\quad
 (L_1,L_2,L_3)=(-8b/3,8b/3,10b/9),
 \qquad \Sigma=\begin{pmatrix}a&b\\b&c\end{pmatrix}.
\]
These explicit linear forms suffice for checking its boundary witnesses.

The coefficient-simplex records reconstruct subdivisions of
$\conv\{0,e_j/\widehat\lambda_j:\widehat\lambda_j>0\}$ by bisecting
a longest edge, with lexicographic tie-breaking. The exact squared
lengths use the coordinates with positive $\widehat\lambda_j$.
Every listed piece must occur in the reconstructed tree and every
terminal piece must be listed. All three current records use matrix
witnesses, with no excluded pieces. The free-coordinate rule in
\cref{cl:sdp-certificate} covers the fourth coordinate in the four-ray
certificate.

The point-rule lower bracket in \cref{cl:counterexample} has a short
explicit witness, independent of the subdivisions:
\[
 \Sigma=\begin{pmatrix}1&-392/621\\-392/621&47753/99290\end{pmatrix},
 \qquad \mu=2539/10000,\qquad
 (\tau_1,\tau_2,\tau_3)=(0,0,227/64).
\]
Its determinant is positive, and direct rational principal-minor checks
give
$\Sigma+\mu\sym(\Sigma N_j)-\tau_jZ\succeq0$ for each ray.
Thus the completed set contains all three points at step $\mu$, proving
$\zBP(\mathbf1)\ge\mu$. The orbit lower bracket $0.97538$ comes from
the rational orbit witness in \cref{fd:counterexample}.

\subsection{A three-cut support-one proof}\label{cl:support-one-certificate}

For \cref{cl:support-one}, the exact ray matrices are
\[
 N_1=\begin{pmatrix}-1&1/4\\0&-3/4\end{pmatrix},\quad
 N_2=\begin{pmatrix}-47/32&1\\-19/32&-3\end{pmatrix},\quad
 N_3=\begin{pmatrix}-25/16&3/8\\-13/16&-9/8\end{pmatrix}.
\]
Take $\Xi_i=\left(\begin{smallmatrix}a_i&b_i+k_i\\b_i-k_i&c_i\end{smallmatrix}\right)$
with the three rows of \cref{cl:three-sets}. They are rows 15, 16, and 59,
numbered from zero, of the retained sixty-set record. The independent
proof below needs only these three rows.

\begin{table}[htbp]
\centering\small
\caption{Rational orbit parameters for the support-one closure lower bound.}
\label{cl:three-sets}
\begin{tabular}{ccccc}\toprule
$i$ &$a_i$ &$b_i$ &$c_i$ &$k_i$\\\midrule
1 &$7753868/7754071$ &$46303/9049648$ &$203/7754071$ &$-9009394/6985127$\\
2 &$1381221/3672875$ &$1434795/2962033$ &$2291654/3672875$ &$-4127451/3621740$\\
3 &$3693704/9822089$ &$3536863/7301600$ &$6128385/9822089$ &$-10687459/9385356$\\
\bottomrule\end{tabular}
\end{table}

The certified steps are
\[
 (\mu_{ij})=
 \begin{pmatrix}
 5557766/5579491&564958/1365417&15591393/8400091\\
 1658819/1667061&15022732/6553787&5144194/6131869\\
 9857221/9905922&19123064/8369713&6479877/7726690
 \end{pmatrix}.
\]
For $\Sigma_i=\sym(\Xi_i)$, direct rational arithmetic gives
$(\Sigma_i)_{11}\ge1/10$ and $\det\Sigma_i\ge10^{-13}$.
Every matrix
$H_{ij}=\sym(\Xi_i(I+\mu_{ij}N_j))$ has
$(H_{ij})_{11}\ge10^{-3}$ and $\det H_{ij}\ge10^{-7}$, except $H_{22}$
and $H_{32}$, whose determinants are at least $10^{-8}$.
All these inequalities follow by inserting the displayed fractions and
cross-multiplying positive denominators. Hence $\Sigma_i\succ0$,
$\det F_i>0$ for $F_i^T=\Xi_iM_b^{-1}$, and
$\alpha_j(C_{F_i})\ge\mu_{ij}$.

Set $G_{ij}=1/\mu_{ij}$. Its inequalities $G_i\lambda\ge1$ are valid
weakenings of the actual intersection cuts. Every point in the full
closure satisfies them. For
\[
 y=(28554901,56139688,14841527)^T/10^8,
\]
the three components of $\mathbf1-G^Ty$ are
\[
 \begin{aligned}
 d_1&=\frac{219837048696181240067}{9087695197410765423400000000},\\
 d_2&=\frac{268223412629374939}{12446434832390989600000000},\\
 d_3&=\frac{3947513071580764783}{157490760258981370980000000}.
 \end{aligned}
\]
They are positive. Therefore $y\ge0$, $G^Ty\le\mathbf1$, and
\[
 \mathbf1^T\lambda\ge y^TG\lambda\ge\mathbf1^Ty
 =24884029/25000000=0.99536116>0.9953611.
\]
This is a complete rational lower-bound proof. The full sixty-set LP
record additionally has the exact value
\[
 \frac{396503562307491652037015395705021177207626376155404352}
 {398351441316353660966368086274075897893464946699521021}
 \simeq0.995361184077.
\]
It is the value of that finite relaxation, not an exact evaluation of the
infinite-family closure. The upper certificate in
\cref{cl:certificate-list} completes the strict nonexactness assertion.

\subsection{Point-rule loss at the three-ray corner}\label{cl:bp-factor}

At \cref{cl:counterexample},
$N_1=\left(\begin{smallmatrix}-7&7/3\\-6&0\end{smallmatrix}\right)$.
For $\Sigma=\left(\begin{smallmatrix}a&b\\b&c\end{smallmatrix}\right)\succ0$,
the kept vector $e_1$ gives $a_1\ge7+6b/a$.
\Cref{cl:kept-boundary} gives $a_1\ge-6b/a$ whenever that bound is
positive; nonnegativity covers the other values. Thus
\[
 a_1\ge\max\{7+6b/a,-6b/a\}\ge7/2.
\]
The symmetric parameter
\[
 \Sigma_*=
 \begin{pmatrix}1&-7/12\\-7/12&7/18\end{pmatrix},\qquad
 \det\Sigma_*=7/144>0,
\]
satisfies $\sym(\Sigma_*(I+(2/7)N_1))=0$.
Its uncompleted set contains the ray endpoint at step $2/7$; its
completion does too. The universal bound proves exact attainment of
$\inf a_1=7/2$.

To bound the corner value away from zero when ray 1 becomes inexpensive,
use the rational orbit parameter
\[
 \Xi=\begin{pmatrix}a&b+k\\b-k&c\end{pmatrix},\qquad
 (a,b,c,k)=
 \left(\frac{658069}{904895},-\frac{83794}{200103},
 \frac{246826}{904895},-\frac{385931}{893930}\right).
\]
Direct principal-minor checks give
\[
 \sym(\Xi)\succ0,\quad \sym(\Xi N_1)\succeq0,\quad
 \sym(\Xi)+(367/759)\sym(\Xi N_j)\succeq0\quad(j=2,3).
\]
Thus ray 1 is recessive, and both other steps are at least $367/759$.
The valid orbit cut implies $\lambda_2+\lambda_3\ge367/759$ for every
$\lambda\in X$. This proves all bounds in the first part of
\cref{cl:unbounded}.

\subsection{The four-ray corner}\label{cl:wcorner-proof}

For the $\famA$ claim use
\[
 Y_1=\begin{pmatrix}141&-63\\-63&29\end{pmatrix},\quad
 Y_2=\begin{pmatrix}4&67\\67&1322\end{pmatrix},\quad
 Y_3=\begin{pmatrix}31&-132\\-132&565\end{pmatrix},\quad Y_4=0.
\]
Their determinants are $120,799,91$. The matrix
$R=-\sum_jN_jY_j=\left(\begin{smallmatrix}14&18\\18&85\end{smallmatrix}\right)$
has determinant $866$. The matrices $20R-Y_j$, $j=1,2,3$, have first
diagonal entries $139,276,249$ and determinants $53340,18479,40551$.
All are positive definite. \Cref{cl:sdp-certificate}, applied to
$\conv\{0,e_1/20,e_2/20,e_3/20\}+\R_+e_4$, shows that every orbit
cut has $a_1+a_2+a_3>1/20$.
Therefore $(20,20,20,0)\in\mathcal C_{\famA}$.

There is also a qualitative explanation of this orbit obstruction.
No $C_F$ can have both $e_w$ and $e_x$ as recession directions; the same
holds for $e_w$ and $-e_y$. To see this, let
$K'=\{d:\sym(F^TM_0(d))\succeq0\}$.
$S$-freeness implies $d_xd_y\le0$ on $K'$ by considering large steps
from an interior point. Thus a small perturbation of $e_x+\sigma e_w$
in the positive $e_y$ direction is outside $K'$, for every $\sigma>0$;
the two-dimensional cone generated by $e_x,e_w$ would lie in its
boundary. If its defining linear map is injective, $K'$ is an invertible
linear image of the $2\times2$ positive-semidefinite cone, whose boundary
cannot contain a two-dimensional convex cone: two nonproportional
rank-one positive-semidefinite matrices have a positive definite sum.
If the map has a nonzero kernel vector $\ell$, then $\ell$ is a
lineality direction of $K'$. The condition
$(e_x+t\ell)_x(e_x+t\ell)_y\le0$ for both signs of small $t$ forces
$\ell_y=0$. For every interior point $s$, both directions along $\ell$
stay in $C_F$, so $q(s+t\ell)\ge0$ for all real $t$.
Its linear coefficient is $\ell_w-s_y\ell_x$.
Variation of $s_y$ over an open interval forces $\ell_x=\ell_w=0$,
a contradiction. The second ray pair follows by exchanging $x$ and $-y$.
The exact $\famB$ completion below has all three positive recession
directions; its parameter $J^TM_b$ is nonsymmetric and hence outside
$\famBP$.

For $\famB$, direct multiplication gives
\[
 \sym(J^TM(s))=
 \begin{pmatrix}-y&(w-1)/2\\(w-1)/2&x\end{pmatrix}.
\]
Its upward completion is the set in the proof of
\cref{cl:unbounded}. With $r=\sqrt{-xy}$, every point of it satisfies
$q\ge1-2r+r^2=(1-r)^2\ge0$, and its interior has $q>0$.
The first three rays are recessive and the fourth leaves at step two,
giving the exact cut vector $(0,0,0,1/2)$.

It remains to prove the sharp $\famBP$ inequality. For
$\Sigma=\left(\begin{smallmatrix}a&b\\b&c\end{smallmatrix}\right)\succ0$,
write $\Delta=ac-b^2$, $A=a+b$, and $B=b+c$.
Here $m=(1,1)^T/2$. The kept forms and the two ray forms are
\[
 v^TZv=\tfrac12v_1\langle\mathbf1,v\rangle_\Sigma,\quad
 v^T\Sigma N_1v=\tfrac12v_2\langle\mathbf1,v\rangle_\Sigma,\quad
 v^T\Sigma N_2v=\tfrac12v_1\langle(1,-1),v\rangle_\Sigma.
\]
The unrestricted maxima of the negative ray forms divided by
$v^T\Sigma v$ are
\begin{equation}\label{cl:w-G}
 G_1=\frac{\sqrt{a(a+2b+c)/\Delta}-1}{4},\qquad
 G_2=\frac{\sqrt{c(a-2b+c)/\Delta}-1}{4}.
\end{equation}
These formulas follow from the two generalized eigenvalue equations:
each has trace $1/2$, and its determinant is respectively
$-(a+b)^2/(16\Delta)$ or $-(c-b)^2/(16\Delta)$.
The uncompleted coefficients are $G_j$; enlarging to the completion
gives $a_j\le G_j$. If a maximizing vector is kept, its necessary
inequality gives the reverse bound and hence $a_j=G_j$.

An exact symmetry exchanges the two rays. Let
\[
 T(x,y,w)=(-y,-x,w),\quad
 H=\begin{pmatrix}1&1\\1&-1\end{pmatrix},\quad
 \Sigma'=H\Sigma^{-1}H,\quad L=\Sigma H/2.
\]
Writing $A_\Sigma(s)=\sym(\Sigma M_b^{-1}M(s))$, direct multiplication
gives
\begin{equation}\label{cl:w-symmetry}
 A_\Sigma(Ts)=L A_{\Sigma'}(s)L^T.
\end{equation}
This congruence preserves the sign of the kept form, maps maximizing
directions to maximizing directions, and exchanges the completed
coefficients: $a_1(\Sigma')=a_2(\Sigma)$ and
$a_2(\Sigma')=a_1(\Sigma)$.
Up to a positive factor,
\[
 \Sigma'=\begin{pmatrix}a+c-2b&c-a\\c-a&a+c+2b\end{pmatrix}.
\]

First suppose $c\le b$. Then $b\ge c>0$, $a>c$, and the positive
maximizer for ray 1 is kept: for $v=(s,-1)$ its quotient is
\[
 f_1(s)=\frac{As-B}{2(as^2-2bs+c)},
\]
whose positive numerator implies $s>B/A>0$ and thus
$v^TZv=s(As-B)/2>0$. Moreover,
\[
 \frac{a(a+2b+c)}{\Delta}
 \ge\frac{a(a+3c)}{c(a-c)}\ge9,
\]
where the last inequality is $(a-3c)^2\ge0$.
Thus $a_1=G_1\ge1/2$.
If instead $a+b\le0$, the transformed parameter satisfies
$c'-b'=2(a+b)/\Delta\le0$, so the same argument gives $a_2\ge1/2$.
Both cases are strictly above the required sum bound.

We may therefore assume $a+b>0$ and $c>b$.
The transformed parameter has the same two strict inequalities.
If both unrestricted maximizing directions are kept, then
\[
 4(a_1+a_2)+2
 =\frac{\sqrt{a(a+2b+c)}+\sqrt{c(a-2b+c)}}{\sqrt\Delta}
 \ge2\sqrt2.
\]
Indeed, AM--GM and
\[
 ac((a+c)^2-4b^2)-4\Delta^2
 =ac(a-c)^2+4b^2\Delta\ge0
\]
prove the inequality. Equality requires $a=c$, $b=0$.

If at least one maximizing direction is not kept, use
\cref{cl:w-symmetry} to make it ray 1. Its positive maximum occurs at
\[
 s_*=
 \frac{a(b+c)+\sqrt{a\Delta(a+2b+c)}}{a(a+b)}.
\]
This formula follows by differentiating $f_1$.
At that maximum the numerator $As_*-B$ is positive, so failure of
keptness means $s_*<0$. Normalize $a=1$ and put $\beta=-b$.
Then $0<\beta<1$, and squaring the negative-numerator condition gives
\[
 \beta^2<c<\frac{2\beta^2}{1+\beta}.
\]
The equivalence uses the identity
\[
 (c-\beta^2)(1-2\beta+c)-(\beta-c)^2
 =(1-\beta)((1+\beta)c-2\beta^2).
\]
The kept vector $(0,-1)$ gives
$a_1\ge(\beta-c)/(2c)$.
The kept-boundary vector $J\Sigma\mathbf1$ gives
$a_2\ge(\beta-c)/(1-2\beta+c)$.
Both denominators are positive. Each ratio decreases with $c$, so
\[
 a_1+a_2>
 \frac{1-\beta}{4\beta}+\beta
 =\beta+\frac1{4\beta}-\frac14\ge\frac34.
\]
This proves the sharp lower bound in every case without a separate
ray-2 admissibility calculation.

For $\Sigma=I$ both maximizing vectors are kept, and direct eigenvalue
calculation on rays 1, 2, and 4 gives
\[
 a=\left((\sqrt2-1)/4,(\sqrt2-1)/4,0,(\sqrt2+1)/4\right).
\]
For ray 4 the endpoint kernel vector $(1,\sqrt2-1)$ is kept, so upward
completion does not extend its step. Ray 3 is recessive by construction.
This proves attainment, uniqueness up to positive scaling, and the
claimed characterization of $(t,t,0,0)$.
At the vertex, the point-rule rotation has angle $\pi/4$ and satisfies
$R_{\pi/4}M_b=\sqrt2I$, so this is also the untransformed point-rule set.

\subsection{An exact local factor lower bound}\label{cl:local-factor}

At \cref{cl:counterexample}, take
$\rc=(11/20,10^{-6},9/20)$ and $z_0=649/2500$.
Let $t_0=\sbar$ and $t_j=\sbar+(z_0/\omega_j)p_j$.
The tetrahedron they span is the image of
$\{\lambda\ge0:\rc^T\lambda\le z_0\}$.
Each triangular facet has rank-two projection onto $(x,y)$: the four
signed projection determinants, in vertex order $012,013,023,123$, are
\[
 11027808/5,\quad 5475613/562500,\quad
 -190382852/45,\quad -1139162725613/562500.
\]
On such a facet the Hessian of $w-xy$ is indefinite, so a minimum cannot
lie in its relative interior. A minimum cannot lie in the tetrahedron's
interior either, since $\partial q/\partial w=1$. It suffices to minimize
on its six edges. Substituting $(1-t)t_i+tt_j$ gives a univariate
quadratic; its exact minimum on $[0,1]$ is listed below.
\[
\begin{array}{c|rrr}
 ij&01&02&03\\\hline
 \min q&3/16&3/2&3/2\\\hline
 ij&12&13&23\\\hline
 \min q&6480592821/4840070400112&19101/15625&
 1361266990991/4320082399920
\end{array}
\]
Every value is positive. Thus $\zK(\rc)>649/2500$.
The completed-family certificate in \cref{cl:certificate-list} gives
$\widehat\lambda=(2211,8805,2464)/10000\in\mathcal C_{\famB}$ and
\[
 \rc^T\widehat\lambda=464971761/2000000000<0.2324859.
\]
Consequently
$\rho_{\famA}\ge\rho_{\famB}>
(649/2500)/(464971761/2000000000)>1.1166$.

\subsection{Retained numerical closure records}\label{cl:numerical-records}

\Cref{cl:survey} gives the eleven-corner survey, all with unit costs
and recorded corner bound one. The single-cut $\famB$ search used restricted
sampling and local refinement, so its values are heuristic lower bounds
on the supremum. Closure cut generation uses finitely many valid cuts:
its LP value is a lower bound on the full closure value. Dividing that
value by the best pricing value found is a numerical correction, not a
certified upper bound, because the pricing search is incomplete.
Equal values in the table therefore indicate observed agreement only.

\begin{table}[htbp]\centering\small
\caption{Recorded single-cut and closure values; arrows show the observed
change. Only the separately stated rational certificates establish exact
brackets.}\label{cl:survey}
\begin{tabular}{lcllll}\toprule
Example or record & support &$\famA$ &$\famB$ &$\mathcal P$ &$\famBP$\\\midrule
\Cref{cl:counterexample}&2&$.97539\to.97539$&$.97539\to.97539$&$.2218$&$.2540\to.2540$\\
Larger orbit example&2&$.83477\to.83477$&$.83477\to.83477$&$.3431$&$.3431\to.3431$\\
R1&2&$.6224\to.6354$&$.7078\to.7921$&$.3532$&$.4238\to.4479$\\
R2&2&$.4526\to.5065$&$.4526\to.5065$&$.1199$&$.1693\to.1693$\\
\Cref{cl:support-one}&1&$.98385\to.99536$&$.98385\to.99536$&$.3986$&$.4272\to.4272$\\
R3&1&$.99923\to.99967$&---&$.4228$&---\\
R4&1&$.98765\to.99574$&---&$.3358$&---\\
R5&1&$.99311\to.99311$&$.99311\to.99311$&$.7706$&$.7934\to.7934$\\
R6&1&$.99525\to.99525$&---&$.5800$&---\\
R7&1&$.99062\to.99062$&---&$.9422$&---\\
R8&1&$.99829\to.99873$&$.99832\to.99874$&$.3742$&$.3910\to.3910$\\
\bottomrule\end{tabular}\end{table}

The larger orbit example is defined after \cref{cl:counterexample}.
R1--R8 identify numerical data records in the electronic companion, which
provides their vertices, ray matrices and source mapping. They are not
additional analytic examples. The survey omits one near-boundary record
because its recorded single-cut values disagree across implementations.
Observed closure gain
was at most about $0.084$ of the corner bound and about $1.12$ times the
single-cut value. The point-rule families gave matching single and
closure estimates in 17 of 18 comparisons: eleven for $\mathcal P$ and
seven for $\famBP$, with record R1 in $\famBP$ the exception.
These observations do not prove convexity of their coefficient up-sets.

For \cref{cl:counterexample} in direction $\mathbf1$, the recorded $\famBP$ values are
close to $16/63$, while the untransformed point-rule set gives about
$0.0896$. Equality with $16/63$ is not established here.
A scan of 291 grid directions and 297 edge directions at this corner
gave largest single-cut loss about $1.1283$ near $(.55,0,.45)$.
At positive costs $(.55,10^{-6},.45)$, the recorded closure estimates for
both $\famA$ and $\famB$ were about $.2301778$, versus corner value
$.2596998$. Only the lower factor bound in \cref{cl:local-factor} is
certified; a finite scan does not upper-bound the worst loss.

The six recorded reoptimization trajectories are
\[
\begin{array}{l|l|r}
 \text{corner}&\text{successive LP bounds}&\text{cuts reaching one}\\\hline
 \text{\Cref{cl:counterexample}}&.97539,.98097,1&3\\
 \text{\Cref{cl:support-one}}&.98385,1&2\\
 \text{R8}&.99829,1&2\\
 \text{larger orbit example}&.83477,.87574,.94583,.98023,.98023,1&6\\
 \text{R2}&.45263,.47486,.47981,.71275,.71275,.72135,.72306,1&8\\
 \text{R1}&.6224,\ldots,.91713&\text{not in 12 rounds}
\end{array}
\]
Each round forms a new basis cone containing the current polyhedron and
maps its orbit cut back to the original coordinates, so the resulting
cuts remain valid for the original feasible set. Zero reduced costs can
make a later corner bound zero; positive-cost theorems then do not
apply. Five of the six trajectories reached the bound, which does not
prove finite convergence of rebasing.

The unresolved extensions are whether a support-one restricted closure
can be exact while no single cut is exact, whether the fixed-corner
$\famA$ and $\famB$ factors at \cref{cl:counterexample} are finite, and why the point-rule
single and closure estimates so often agree. The support-one $\famB$
closure upper certificate remains incomplete. No claim here depends on
that incomplete record or on the stopped four-ray numerical search;
the four-ray result has the analytic proof above.
